\chapter{घटाव: पहले कदम}\label{ch:g1:subtraction}

सात गुब्बारे, दो उड़ गए: कितने बचे? कम कर देना ही
\emph{घटाना} है, जिसे $-$ चिह्न से लिखते हैं। घटाव का एक दूसरा
रूप भी है: यह ``कितने ज़्यादा?'' का उत्तर भी देता है।

\section{निकाल लेना}

\begin{definition}[घटाव]\label{def:g1:subtraction:def}
\emph{घटाना} यानी कुछ निकाल लेना और जो बचा उसे गिनना। हम
लिखते हैं
\[
7 - 2 = 5
\]
और पढ़ते हैं: ``सात घटा दो बराबर पाँच''।
\end{definition}

\begin{omfigure}
\begin{tikzpicture}[scale=0.6]
  \foreach \i in {0,...,4} \fill[omDef] (\i*1.05,0) circle (0.3);
  \foreach \i in {5,6} {
    \fill[omDef!30] (\i*1.05,0) circle (0.3);
    \draw[omThm, thick] (\i*1.05-0.32,-0.32) -- (\i*1.05+0.32,0.32);
    \draw[omThm, thick] (\i*1.05-0.32,0.32) -- (\i*1.05+0.32,-0.32);
  }
  \node[right, font=\small] at (7.6,0)
    {$7 - 2 = 5$: दो काटो, पाँच बचे};
\end{tikzpicture}

{\small चित्रों में घटाव: काट देना ही निकाल लेना है।}
\end{omfigure}

\begin{example}[पीछे गिनना]\label{ex:g1:subtraction:countback}
थोड़ा-सा घटाना हो तो रेखा पर पीछे गिनो: $11 - 3$: कहो
``ग्यारह'', फिर तीन कदम पीछे --- ``दस, नौ, आठ''। तो
$11 - 3 = 8$।
\end{example}

\section{कितने ज़्यादा?}

\begin{example}[अंतर]\label{ex:g1:subtraction:difference}
``नीना के पास $9$ कंचे हैं और टॉम के पास $6$। नीना के पास कितने ज़्यादा?''
यहाँ कुछ निकाला नहीं गया, फिर भी यह घटाव है: कंचों के जोड़े बनाओ
और नीना के वे कंचे गिनो जिनका जोड़ीदार नहीं:
\[
9 - 6 = 3 .
\]
नीना के पास $3$ ज़्यादा हैं। घटाव के उत्तर को
\emph{अंतर}\index{अंतर} कहते हैं।
\end{example}

\begin{method}[आगे गिनना]\label{met:g1:subtraction:countup}
\omterm{ex:g1:subtraction:difference}{अंतर} निकालने के लिए \emph{आगे} गिनना अक्सर सबसे आसान होता है:
$12 - 9$ के लिए $9$ से शुरू करके $12$ तक चढ़ो: ``दस, ग्यारह,
बारह'' --- तीन कदम। तो $12 - 9 = 3$। (बड़े फ़ासले के लिए $10$ से
होकर चढ़ो: $14 - 8$: $8$ से $10$ तक $2$, $10$ से $14$
तक $4$: \omterm{ex:g1:subtraction:difference}{अंतर} $6$।)
\end{method}

\begin{omfigure}
\begin{tikzpicture}[scale=0.6]
  \draw[-Stealth] (5.6,0) -- (15.4,0);
  \foreach \i in {6,...,15} \draw (\i,-0.1) -- (\i,0.1);
  \foreach \i in {8,10,14}
    \node[below, font=\footnotesize] at (\i,-0.12) {\i};
  \draw[-Stealth, omDef, thick] (8,0.3) .. controls (9,0.9) ..
    (10,0.3);
  \node[omDef, font=\small] at (9,1.0) {$+2$};
  \draw[-Stealth, omThm, thick] (10,0.3) .. controls (12,1.2) ..
    (14,0.3);
  \node[omThm, font=\small] at (12,1.25) {$+4$};
\end{tikzpicture}

{\small आगे गिनकर $14 - 8$: $8$ से $14$ तक
$2 + 4 = 6$ कदम हैं।}
\end{omfigure}

\section{जोड़ और घटाव जुड़वाँ हैं}

\begin{example}[तथ्यों का परिवार]\label{ex:g1:subtraction:family}
संख्याएँ $3$, $4$ और $7$ मिलकर एक छोटा परिवार बनाती हैं:
\[
3 + 4 = 7, \qquad
4 + 3 = 7, \qquad
7 - 4 = 3, \qquad
7 - 3 = 4 .
\]
परिवार का एक तथ्य पता हो तो बाक़ी तीन मुफ़्त में मिल जाते हैं।
घटाव जोड़ को उलट देता है: $+4$ के बाद $-4$ करो तो वापस वहीं।
\end{example}

\begin{example}[जाँच]\label{ex:g1:subtraction:check}
$13 - 5 = 8$ निकाला? जुड़वाँ जोड़ से जाँचो:
$8 + 5 = 13$। \checkmark\ घटाव ठीक तभी है जब वापस जोड़ने पर
पहली संख्या मिल जाए।
\end{example}

\begin{example}[कौन-सी संक्रिया?]\label{ex:g1:subtraction:choose}
``तार पर $6$ चिड़ियाँ, $4$ उड़ जाती हैं'': निकालना, $6 - 4 = 2$।
``$6$ चिड़ियाँ, फिर $4$ और आ जाती हैं'': मिलाना, $6 + 4 = 10$।
``$6$ गौरैया और $4$ रॉबिन: गौरैया कितनी ज़्यादा?'':
\omterm{ex:g1:subtraction:difference}{अंतर}, $6 - 4 = 2$। पक्का न लगे तो कहानी का चित्र बनाओ!
\end{example}

\section{अभ्यास}

\begin{exercise}[$\star$]\label{exo:g1:subtraction:1}
बनाओ, काटो और निकालो: $6 - 2$; \quad $8 - 5$; \quad
$10 - 4$।
\end{exercise}

\begin{exercise}[$\star$]\label{exo:g1:subtraction:2}
पीछे गिनकर निकालो: $9 - 3$; \quad $12 - 2$; \quad
$15 - 4$।
\end{exercise}

\begin{exercise}[$\star$]\label{exo:g1:subtraction:3}
आगे गिनकर \omterm{ex:g1:subtraction:difference}{अंतर} निकालो: $11 - 8$; \quad $13 - 9$; \quad
$16 - 12$।
\end{exercise}

\begin{exercise}[$\star$]\label{exo:g1:subtraction:4}
$2$, $6$ और $8$ के परिवार के चारों तथ्य लिखो।
\end{exercise}

\begin{exercise}[$\star$]\label{exo:g1:subtraction:5}
निकालो, फिर जुड़वाँ जोड़ से जाँचो: $14 - 6$; \quad
$17 - 9$।
\end{exercise}

\begin{exercise}[$\star$]\label{exo:g1:subtraction:6}
ख़ाली जगह भरो:
\[
10 - \;?\; = 7, \qquad
\;?\; - 5 = 5, \qquad
12 - \;?\; = 12 .
\]
\end{exercise}

\begin{exercise}[$\star$]\label{exo:g1:subtraction:7}
``कटोरे में $15$ चेरी हैं। बच्चे $8$ खा जाते हैं। कितनी बचीं?''
घटाव और जवाब का वाक्य लिखो।
\end{exercise}

\begin{exercise}[$\star$]\label{exo:g1:subtraction:8}
हर कहानी के लिए $+$ या $-$ चुनो, फिर निकालो: ``तालाब पर $9$
बत्तखें, $5$ तैरकर चली जाती हैं''; ``$9$ बत्तखें, $5$ हंस आ जाते हैं'';
``अली $9$ साल का है, उसकी बहन $5$ की: अली कितने साल बड़ा है?''।
\end{exercise}

\begin{exercise}[$\star\star$]\label{exo:g1:subtraction:9}
बत्तख के पास $12$ अंडे थे। कुछ से बच्चे निकल आए; अब $12 - ?$ अंडे
बचे हैं और वह $7$ बिना फूटे अंडे गिनती है। कितने अंडों से बच्चे
निकले? (यह किस तथ्य-परिवार का सवाल है?)
\end{exercise}
